그래프(Graph)는 정보 문제에서 큰 비중을 차지하는 분야로 정점과 정점을 연결하는 간선으로 이루어진 집합의 일종이다. 이번 챕터에서는 이 그래프의 정의외 그래프 자료의 저장 방식, 그리고 중로 사용되는 몇 가지 용어를 알아본다.

\section{그래프의 정의}
그래프의 엄밀한 정의는 이와 다르지만, PS 공부에 있어서는 다음과 같이 간략하게 이해하여도 충분하다.

\begin{definition}
    그래프 $G$는 정점의 집합 $V$와 간선의 집합 $E$에 대해 $(V, E)$의 순서쌍을 나타낸다. 또한, $E$의 임의의 원소 $e$는 $V$에 속하는 임의의 두 원소 $u$와 $v$로 이루어진 순서쌍이어야 한다.
    이러한 $V$와 $E$에 대해 그래프 $G$는 $V$에 속하는 정점으로 이루어져 있고, $E$에 속하는 임의의 원소 $e$의 성분 두 정점이 연결된 것을 나타낸다.
\end{definition}

일반적으로 그래프는 어떠한 두 개체의 연결관계를 표시할 때 유용하다. 많은 그래프 문제에서 직접 그래프를 제공하기도 하지만, 떄때로는 주어진 상황으로부터 직접 그래프를 고안해야 함에 유의하라.
대표적인 예시로는 격자 사이에서 인접한 칸만 이동할 수 있는 경우를 그래프로 나타낼 때가 있을 것이다. $N \times M$ 격자를 그래프로 나타내고 그 이동을 시뮬레이션 할 때, $V$를 모든 격자의 각 칸을 모은 집합, $E$를 각 칸에 대해 인접한 네 칸과 연결한 간선으로 정의하면 그래프로 표현할 수 있게 된다.

\section{그래프의 표현}
우리는 그래프 사진을 보고 직관적으로 그래프의 구조와 모습을 이해할 수 있지만, 컴퓨터는 그렇지 않다. 우리는 직접 그래프의 모양을 적절한 규칙에 따라 입력해주어야 한다.
일반적으로 정점은 숫자 하나로 나타낸다. 정점의 개수가 $N$개인 정점의 경우에는 $0$부터 $N-1$까지 또는 $1$부터 $N$까지를 사용한다.
일반적으로 간선은 총 $M$개의 줄에 걸쳐 각 줄마다 써진 $u$와 $v$를 입력받는데, 다른 말로 하면 $E$의 원소를 직접 입력받는다는 의미이다.

하지만 구현의 관점에서는 정점 $E$의 집합을 입력받는 방식은 유용하지 않은데, 만약 어떤 정점 $v$에서 갈 수 있는 모든 정점을 탐색하려면 $O(M)$의 시간 복잡도가 필요하기 때문이다. 
따라서 일반적으로 인접행렬 또는 인접리스트를 사용하여 그래프를 표시한다.

\begin{itemize}
    \item 인접행렬은 $i$행 $j$열의 성분이 정점 $i$에서 정점 $j$로 갈 수 있는지 여부를 저장한 행렬이다. $O(N^2)$의 공간과 $v$에서 갈 수 있는 모든 정점 탐색에 $O(N)$의 시간을 요구한다.
    \item 인접리스트는 각 정점 $i$에 대해 $i$에서 갈 수 있는 모든 정점을 저장한 리스트이다. 보통 \texttt{std::vector}를 사용해 구현하며, $O(M)$의 공간과 $v$에서 갈 수 있는 모든 정점 탐색에 연결된 정점 개수만큼의 시간을 요구한다. 다른 말로 낭비되는 반복이 없다.
\end{itemize}

다음은 인접행렬과 인접리스트를 만드는 예시 코드이다.
\begin{code} \label{code: adj-matrix}
인접행렬 구현 코드
\begin{lstlisting}
#include <bits/stdc++.h>
using namespace std;

int n, m;
int adj[1010][1010];

int main() {
    ios_base::sync_with_stdio(false);
    cin.tie(0);
    cin >> n >> m;
    for (int i = 1; i <= m; i++) {
        int u, v;
        cin >> u >> v;
        adj[u][v] = 1;
        adj[v][u] = 1;
    }
}
\end{lstlisting}
\end{code}

\begin{code} \label{code: adj-list}
인접리스트 구현 코드
\begin{lstlisting}
#include <bits/stdc++.h>
using namespace std;

int n, m;
vector<int> adj[101010];

int main() {
    ios_base::sync_with_stdio(false);
    cin.tie(0);
    cin >> n >> m;
    for (int i = 1; i <= m; i++) {
        int u, v;
        cin >> u >> v;
        adj[u].push_back(v);
        adj[v].push_back(u);
    }
}
\end{lstlisting}
\end{code}

\section{몇 가지 용어들}
그래프 이론에서 사용되는 몇 가지 용어가 있다. 다음 용어들은 이 파트에서 꾸준이 사용될 것이므로 꼭 숙지하도록 하자.

\begin{itemize}
    \item 연결: 두 정점 사이를 간선을 통해 이동할 수 있음을 나타내는 용어
    \item 경로: 연결된 두 정점 사이를 간선을 통해 이동할 때 지나가는 간선의 나열
    \item 최단경로: 어떤 두 정점 사이를 있는 모든 경로에 대해 경로의 길이가 최소인 경로
    \item 사이클: 한 정점에서 출발해 다시 그 정점으로 돌아오는 경로 중 경로를 이루는 간선이 모두 다른 경로
    \item 차수: 한 정점과 연결된 다른 정점의 수. 다른 말로 인접 리스트에서 리스트의 길이
    \item 컴포넌트: 임의의 두 정점 사이를 간선을 통해 이동할 수 있는 $V$의 부분집합
    \item 완전 그래프: 각 정점 사이를 정확히 하나의 간선으로 연결하는 그래프 ($K_n$)
    \item 클릭: 임의의 두 정점 사이가 간선을 통해 연결되어 있는 $G$의 부분집합, 즉 부분 완전 그래프
    \item 
\end{itemize}

\section{연습문제 A}
\begin{problem}
    \Cref{code: adj-matrix}와 \Cref{code: adj-list}를 보고 어떤 그래프 하나의 입력 데이터를 직접 만들어보라.
\end{problem}

\section{연습문제 B}

\subsection{기초문제}
\begin{problem} \url{https://www.acmicpc.net/problem/3132} \end{problem}
\begin{problem} \url{https://www.acmicpc.net/problem/5247} \end{problem}
\begin{problem} \url{https://www.acmicpc.net/problem/9993} \end{problem}
\begin{problem} \url{https://www.acmicpc.net/problem/10328} \end{problem}
\begin{problem} \url{https://www.acmicpc.net/problem/10965} \end{problem}
\begin{problem} \url{https://www.acmicpc.net/problem/11255} \end{problem}
\begin{problem} \url{https://www.acmicpc.net/problem/14676} \end{problem}
\begin{problem} \url{https://www.acmicpc.net/problem/17511} \end{problem}
\begin{problem} \url{https://www.acmicpc.net/problem/18097} \end{problem}
\begin{problem} \url{https://www.acmicpc.net/problem/18788} \end{problem}
\begin{problem} \url{https://www.acmicpc.net/problem/22035} \end{problem}
\begin{problem} \url{https://www.acmicpc.net/problem/25953} \end{problem}
\begin{problem} \url{https://www.acmicpc.net/problem/29955} \end{problem}

\subsection{응용문제}
\begin{problem} \url{https://www.acmicpc.net/problem/1050} \end{problem}
\begin{problem} \url{https://www.acmicpc.net/problem/1533} \end{problem}
\begin{problem} \url{https://www.acmicpc.net/problem/2008} \end{problem}
\begin{problem} \url{https://www.acmicpc.net/problem/2099} \end{problem}
\begin{problem} \url{https://www.acmicpc.net/problem/2549} \end{problem}
\begin{problem} \url{https://www.acmicpc.net/problem/3378} \end{problem}
\begin{problem} \url{https://www.acmicpc.net/problem/4165} \end{problem}
\begin{problem} \url{https://www.acmicpc.net/problem/6691} \end{problem}
\begin{problem} \url{https://www.acmicpc.net/problem/11642} \end{problem}
\begin{problem} \url{https://www.acmicpc.net/problem/11749} \end{problem}
\begin{problem} \url{https://www.acmicpc.net/problem/13143} \end{problem}
\begin{problem} \url{https://www.acmicpc.net/problem/19544} \end{problem}
\begin{problem} \url{https://www.acmicpc.net/problem/19934} \end{problem}

\subsection{도전문제}
\begin{problem} \url{https://www.acmicpc.net/problem/2939} \end{problem}
\begin{problem} \url{https://www.acmicpc.net/problem/11755} \end{problem}
\begin{problem} \url{https://www.acmicpc.net/problem/15261} \end{problem}
\begin{problem} \url{https://www.acmicpc.net/problem/16661} \end{problem}
\begin{problem} \url{https://www.acmicpc.net/problem/18346} \end{problem}
\begin{problem} \url{https://www.acmicpc.net/problem/18971} \end{problem}
\begin{problem} \url{https://www.acmicpc.net/problem/19221} \end{problem}
\begin{problem} \url{https://www.acmicpc.net/problem/20689} \end{problem}
\begin{problem} \url{https://www.acmicpc.net/problem/21117} \end{problem}
\begin{problem} \url{https://www.acmicpc.net/problem/22025} \end{problem}
\begin{problem} \url{https://www.acmicpc.net/problem/26134} \end{problem}
\begin{problem} \url{https://www.acmicpc.net/problem/27167} \end{problem}
\begin{problem} \url{https://www.acmicpc.net/problem/29257} \end{problem}
\begin{problem} \url{https://www.acmicpc.net/problem/30042} \end{problem}